\section{本讲小结}

德布罗意波长由指定参考系内的动量决定。电场加速提供动能，非相对论时以 $p=\sqrt{2mK}$ 求波长，高速时须使用 $p^2c^2=K(K+2mc^2)$。物质波描述概率幅的相位与叠加，探测到的单个粒子仍形成局域记录，波线不是粒子的经典轨迹。

玻尔模型将圆周运动条件与角动量量子化联立，得到类氢半径、速率和能级。内部总能量差决定跃迁能标，电离极限规定束缚能的零点；光谱计数须区别单个原子的下降路径与大量原子的所有可能路径。类氢公式要求单电子体系，并须检查固定核、低速等近似。

有限核质量通过约化质量修正内部能级；原子整体反冲则由发射、吸收过程的能量和动量守恒计算。这两种质量效应分别属于内部相对运动和整体运动。自由静止原子的发射光子能量低于内部能级差，吸收所需能量高于该差；微小修正应先保留精确关系，再按小量展开。

粒子碰撞激发时，可供内部状态改变的是质心系的相对运动动能。实验室的整体平动由总动量约束，不能任意转为激发能。静靶、相向运动以及不同入射质量会给出不同的实验室阈能；达到阈值仅表示能量、动量允许，实际激发和随后谱线的出现还取决于碰撞与跃迁过程。
